%% Chapter 2 — The Breed Category and Formal Signatures

\chapter{The Breed Category and Formal Signatures}
\label{chap:breed-category}

\epigraph{%
  ``A category is a collection of objects and morphisms satisfying identity and
  associativity.''
}{--- Saunders Mac Lane, \textit{Categories for the Working Mathematician}, 1971}

\section{Breed Signatures}
\label{sec:breed-signatures}

The central organizing concept of the Periodic Table of Reason is the
\emph{breed}: a formal specification of an inference modality together with a
concrete implementation that is verifiable, falsifiable, and receipt-bound.
Before defining the breed category, we fix the algebraic signature that every
breed must possess.

\begin{definition}[Breed Signature]
\label{def:breed-signature}
A \emph{breed signature} is a triple $\Sigma_i = (X_i, Y_i, S_i)$ where:
\begin{itemize}
  \item $X_i$ is the \emph{input type}: a measurable space of admissible inputs to
    breed $i$, equipped with the $\sigma$-algebra generated by the JSON schema
    declared in the breed's contract.
  \item $Y_i$ is the \emph{output type}: a measurable space of admissible outputs,
    whose elements are records conforming to the \texttt{ContractResult} shape
    \[
      Y_i \ni y = \bigl(\,\text{status},\; \Breed_i,\; \text{run\_id},\;
                         \text{output\_hash},\; \text{replay\_pointer},\;
                         \text{options\_profile},\; \text{input\_hash},\;
                         \text{output}\bigr).
    \]
  \item $S_i \subseteq X_i \times Y_i$ is the \emph{specification relation}:
    the set of $(x, y)$ pairs that constitute correct behavior for breed $i$.
\end{itemize}
\end{definition}

\begin{remark}
The field \texttt{input\_hash} in $Y_i$ is not redundant with the external receipt
chain.  It is computed \emph{inside} the WASM module—specifically in
\texttt{cognition\_run()} at \texttt{crates/wasm4pm-cognition/src/wasm.rs}—before
any output is generated.  Its inclusion in $Y_i$ means that every element of the
output type carries a self-certifying reference to the input that caused it.  This
makes the output type \emph{causally complete}: the pair $(x, y)$ can be verified
as a valid member of $S_i$ without any external context, because $y$ contains
$\blake(x)$ as a first-class field.  The addition of \texttt{input\_hash} is
additive (no existing field is removed or retyped) and strictly strengthens the
causal completeness property of $Y_i$.
\end{remark}

\begin{definition}[Breed Implementation]
\label{def:breed-implementation}
Given a breed signature $\Sigma_i = (X_i, Y_i, S_i)$, a \emph{breed
implementation} is a measurable function $\hat{B}_i : X_i \to Y_i$ such that
\[
  \forall x \in X_i,\quad (x,\, \hat{B}_i(x)) \in S_i.
\]
The canonical implementation is the WebAssembly module compiled from the breed's
Rust source in \texttt{crates/wasm4pm-cognition/}.
\end{definition}

\section{The Breed Category \texorpdfstring{$\mathbf{Breed}$}{Breed}}
\label{sec:breed-category}

We now assemble the breed signatures into a category.

\begin{definition}[The Breed Category]
\label{def:breed-category}
The \emph{breed category} $\mathbf{Breed}$ is defined as follows.
\begin{itemize}
  \item \textbf{Objects.}  The objects of $\mathbf{Breed}$ are the breed signatures
    $\Sigma_i = (X_i, Y_i, S_i)$ for $i$ ranging over the thirteen breeds
    registered in the current platform (see Chapter~\ref{chap:periodic-table}).
  \item \textbf{Morphisms.}  A morphism $\phi : \Sigma_i \to \Sigma_j$ is a pair of
    measurable maps $(\phi_X : X_j \to X_i,\; \phi_Y : Y_i \to Y_j)$---note the
    contravariance in the input direction---such that
    \[
      (x, y) \in S_i \implies (\phi_X^{-1}(x),\; \phi_Y(y)) \in S_j.
    \]
    Such a morphism is called a \emph{behavior-preserving transformation}: it
    translates inputs from $j$'s domain into $i$'s domain, runs $i$'s specification,
    and translates the output into $j$'s output type, with the guarantee that correct
    behavior is preserved.
  \item \textbf{Identity.}  For each $\Sigma_i$, the identity morphism
    $\mathrm{id}_{\Sigma_i}$ is the pair $(\mathrm{id}_{X_i}, \mathrm{id}_{Y_i})$.
  \item \textbf{Composition.}  Given $\phi : \Sigma_i \to \Sigma_j$ and
    $\psi : \Sigma_j \to \Sigma_k$, the composite $\psi \circ \phi : \Sigma_i \to
    \Sigma_k$ is defined by
    \[
      (\psi \circ \phi)_X = \phi_X \circ \psi_X, \qquad
      (\psi \circ \phi)_Y = \psi_Y \circ \phi_Y.
    \]
\end{itemize}
\end{definition}

\begin{proposition}[Each Breed as a Functor]
\label{prop:breed-functor}
For each breed $i$, the canonical implementation $\hat{B}_i$ extends to a functor
\[
  \mathcal{F}_i : \mathbf{Meas}_{/X_i} \longrightarrow \mathbf{Meas}_{/Y_i}
\]
from the slice category of measurable spaces over $X_i$ to the slice category over
$Y_i$, defined on objects by $\mathcal{F}_i(A \xrightarrow{f} X_i) = Y_i
\xrightarrow{\hat{B}_i} Y_i$ and on morphisms by post-composition with $\hat{B}_i$.
\end{proposition}

\begin{proof}
We verify the two functor laws.

\medskip
\noindent\textbf{Identity preservation.}  Let $A \xrightarrow{f} X_i$ be an object
of $\mathbf{Meas}_{/X_i}$ and let $\mathrm{id}_f$ be its identity morphism.  By
definition, $\mathcal{F}_i(\mathrm{id}_f) = \mathrm{id}_{\hat{B}_i \circ f}$, which
is the identity morphism on $\mathcal{F}_i(A \xrightarrow{f} X_i)$.  Thus
$\mathcal{F}_i$ preserves identities.

\medskip
\noindent\textbf{Composition preservation.}  Let $h : (A \xrightarrow{f} X_i) \to
(B \xrightarrow{g} X_i)$ and $k : (B \xrightarrow{g} X_i) \to (C \xrightarrow{m}
X_i)$ be morphisms in $\mathbf{Meas}_{/X_i}$, so that $g \circ h = f$ and $m \circ
k = g$ as measurable maps.  Then $\mathcal{F}_i(k \circ h)$ is defined by
post-composition: $\hat{B}_i \circ (m \circ k \circ h) = \hat{B}_i \circ m \circ k
\circ h$.  On the other hand, $\mathcal{F}_i(k) \circ \mathcal{F}_i(h) =
(\hat{B}_i \circ m \circ k) \circ (\hat{B}_i \circ g \circ h)$.  Since $\hat{B}_i$
is a fixed measurable function (not a natural transformation between two different
functors), the post-composition operation is strictly associative, so both
expressions reduce to $\hat{B}_i \circ m \circ k \circ h$.  Therefore
$\mathcal{F}_i(k \circ h) = \mathcal{F}_i(k) \circ \mathcal{F}_i(h)$, and
composition is preserved.

\medskip
\noindent The functor $\mathcal{F}_i$ is well-defined because $\hat{B}_i$ is
measurable (being a composition of measurable maps: the JSON schema decoder,
the WASM execution, and the BLAKE3 hash), so post-composition with $\hat{B}_i$
sends measurable maps to measurable maps.  This completes the proof.
\end{proof}

\section{The Implementation Space}
\label{sec:implementation-space}

Definition~\ref{def:breed-implementation} identifies one canonical implementation
$\hat{B}_i$ for each breed.  In the proof of operational falsifiability, however,
we must reason about the \emph{space of all possible implementations}—including
non-conforming ones—in order to show that the test suite separates the canonical
implementation from every alternative.

\begin{definition}[Implementation Space]
\label{def:implementation-space}
Given a breed signature $\Sigma_i = (X_i, Y_i, S_i)$, the \emph{implementation
space} is the function space
\[
  Y_i^{X_i} = \{ f : X_i \to Y_i \mid f \text{ is measurable} \},
\]
equipped with the compact-open topology: for each compact $K \subseteq X_i$ and
each open $U \subseteq Y_i$, the set $V(K, U) = \{ f \in Y_i^{X_i} \mid f(K)
\subseteq U \}$ is declared open.
\end{definition}

\begin{remark}
For finite $X_i$, every subset of $X_i$ is compact and the compact-open topology
on $Y_i^{X_i}$ coincides with the product topology, which on a finite product of
discrete spaces is itself discrete.  In the discrete topology, every singleton
$\{f\}$ is open, so the specification relation $S_i$ (viewed as a subset of
$Y_i^{X_i}$ via the identification $f \leftrightarrow \{(x, f(x))\}_{x \in X_i}$)
is a singleton $\{ \hat{B}_i \}$.  This means that for finite input spaces, the
specification \emph{uniquely determines} the implementation: there is no
topological neighborhood around $\hat{B}_i$ that contains a non-conforming
alternative, and the notion of ``approximation'' collapses.  The interesting case
is therefore the infinite or continuous $X_i$, where $S_i$ can be a proper closed
subset of $Y_i^{X_i}$ and the compact-open topology governs the sense in which
one implementation can be ``close to'' another without being identical.
\end{remark}

\begin{proposition}[Specification as Closed Subspace]
\label{prop:spec-closed}
Under the compact-open topology on $Y_i^{X_i}$, the set of conforming
implementations $\mathcal{C}_i = \{ f \in Y_i^{X_i} \mid \forall x,\, (x,f(x)) \in
S_i \}$ is a closed subspace whenever $S_i$ is a closed relation in $X_i \times
Y_i$.
\end{proposition}

\begin{proof}
The evaluation map $\mathrm{ev} : Y_i^{X_i} \times X_i \to Y_i$ defined by
$\mathrm{ev}(f, x) = f(x)$ is continuous in the compact-open topology when $X_i$
is compactly generated \cite{maclane1971}.  The conformance condition
$(x, \mathrm{ev}(f,x)) \in S_i$ defines $\mathcal{C}_i$ as the preimage of $S_i$
under the continuous map $f \mapsto \{(x, f(x))\}_{x \in X_i}$.  Since the
preimage of a closed set under a continuous map is closed, $\mathcal{C}_i$ is
closed.

More concretely: suppose $(f_\alpha)$ is a net in $\mathcal{C}_i$ converging to
some $f_\infty \in Y_i^{X_i}$ in the compact-open topology.  Convergence in the
compact-open topology implies pointwise convergence on compact sets; since every
singleton $\{x\}$ is compact, $f_\alpha(x) \to f_\infty(x)$ for all $x \in X_i$.
Because $S_i$ is closed in $X_i \times Y_i$, and $(x, f_\alpha(x)) \in S_i$ for
all $\alpha$, the limit $(x, f_\infty(x)) \in S_i$ as well.  Therefore $f_\infty
\in \mathcal{C}_i$, confirming that $\mathcal{C}_i$ is closed.
\end{proof}

\section{Summary}
\label{sec:ch2-summary}

This chapter formalized the breed signature $(X_i, Y_i, S_i)$ and assembled the
thirteen breed signatures into the category $\mathbf{Breed}$, whose morphisms are
behavior-preserving transformations.  The canonical implementation of each breed
was shown to extend to a functor on slice categories, and the implementation space
was given the compact-open topology in which the conforming implementations form a
closed subspace.  Chapter~3 introduces the Certainty Factor calculus of the MYCIN
breed as the first concrete occupant of the Periodic Table, grounding the abstract
definitions of this chapter in a historically significant and fully receipt-bound
inference system.
